\documentstyle[%


righttag%


,11pt%


,amssymb%


,russian%


,izvru%


]{amsart}





\newcommand{\tl}[3]{\medskip\noindent{\bf#1}\ #2 \dotfill #3 \par} 


\pagestyle{empty}


\begin{document}


%\tableofcontents


\par


\vspace*{1cm}


\par


\bigskip





Данный тематический номер журнала является продолжением предыдущего


номера~7 за 2004 год, посвященного памяти профессора


КГУ Чибриковой Л.И.





\bigskip


\bigskip


\par


\centerline{\Large СОДЕРЖАНИЕ}


\bigskip


\bigskip


%%%%%%%%%% Use \newline %%%%%%%%%%%














\tl{Агачев Ю.Р.}{Об оптимизации прямых  методов решения обыкновенных  


интегродиффе-\newline ренциальных уравнений}{3}


\tl{Габдулхаев Б.Г.}


{Методы решения бисингулярных интегральных уравнений


с внутрен-\newline ними коэффициентами}{11}


\tl{Дюкарев Ю.М.}


{О неопределенности интерполяционных задач


для неванлинновских\newline функций}{26}


\tl{Ермолаева Л.Б.}


{Решение периодическиx краевых задач методом подобластей}{39}


\tl{Насыров С.Р.}


{Римановы поверхности, ограниченные кривыми,


с заданными проекциями\newline точек ветвления}{48}


\tl{Сизиков В.С., Смирнов А.В., Федоров Б.А.}


{Численное решение сингулярного инте-\newline грального уравнения Абеля


обобщенным методом квадратур}{62}


\tl{Тихоненко Н.Я.}{Конечномерные методы приближенного


решения линейных уравнений\newline типа свертки}{71}














\vfill


\noindent\raisebox{2pt}{\copyright} Казанский госудаpственный унивеpситет


\end{document}


